\section{Related Work}

Hsiao's odd-weight-column construction makes checker complexity a code-design concern~\cite{hsiao1970}. Ghosh \emph{et al.} optimized Hamming/Hsiao checkers for power while monitoring area and delay~\cite{ghosh2004}, and Reviriego \emph{et al.} constructed low-delay SEC/SECDED logic~\cite{reviriego2013}. Recent ASIC synthesis compares SECDED and SECDED-DAEC designs over several lengths~\cite{tripathi2020}. These works establish hardware-conscious SECDED; our contribution is not a lower-cost code, but a matched post-route test of architecture identity after semantic equivalence.

Naseer and Draper compared parallel SEC, Hsiao SECDED, DEC, and DEC-TED hardware~\cite{naseer2008}; Layered-ECC targets low-complexity DEC decoding~\cite{das2019}. Wong \emph{et al.} altered Chien-search power management~\cite{wong2011}, while recent direct and conventional BCH decoders show that architecture can materially change area efficiency and latency~\cite{lagendijk2026}. This prior art requires our BCH conclusion to remain tied to the evaluated syndrome/Chien RTL.

STT-MRAM studies co-design correction strength with device/array cost~\cite{delbel2014,pajouhi2015}; LAD-ECC and Stealth ECC adapt protection using workload or data properties~\cite{wei2020,lee2022}; FaultSim supplies field-model-oriented system reliability evaluation~\cite{nair2016}. We instead hold protection semantics and useful work fixed, measure codec logic in a pinned open flow~\cite{ajayi2019}, and make no FIT, SER, array, or operational-failure claim. A focused 2018--2026 audit found no close paper combining exact RTL identity transfer, common post-route constraints, predeclared matched seeds, feasibility-gated energy, and explicit unavailable points. We use that only to position the experiment; it is not a priority or universal-absence claim.
